\documentclass[11pt,a4paper,oneside]{book}
\usepackage[margin=2.5cm]{geometry}
\usepackage{amsmath}
\usepackage{booktabs}
\usepackage{verbatim}
\usepackage{xcolor}
\usepackage[colorlinks=true,linkcolor=blue!50!black,urlcolor=blue!60!black]{hyperref}

\newcommand{\key}[1]{\texttt{#1}}
\newcommand{\file}[1]{\texttt{#1}}
\newenvironment{shell}{\small\verbatim}{\endverbatim}
\newcommand{\deck}[1]{{\small\verbatiminput{#1}}}

\title{\Huge VESPAW\\[0.5em]\Large Tutorial\\[2em]
\normalsize Hands-on examples with runnable, checked input decks}
\author{Mohsen Sotoudeh}
\date{Version of \today}

\begin{document}
\frontmatter
\maketitle
\tableofcontents
\mainmatter

% =====================================================================
\chapter*{How to use this tutorial}
\addcontentsline{toc}{chapter}{How to use this tutorial}
% =====================================================================

Every example in this tutorial lives in a directory
\file{examples/tutorial/NN\_name/} and is run by
\begin{shell}
python3 doc/src/tutorial/check_examples.py             # all examples
python3 doc/src/tutorial/check_examples.py 02_si       # one example
python3 doc/src/tutorial/check_examples.py 02_si:scf   # one run of it
\end{shell}
The input files in these directories are \emph{templates}: words such
as \texttt{@PSIDYN@} or \texttt{@LUNIT@} are placeholders. The file
\file{case.json} lists the runs of an example and what each placeholder
becomes in each run. The checker writes the actual inputs into
\file{examples/tutorial/NN\_name/work/RUN/} and runs them there. After a
run you find in that directory exactly the input that was used, the report
(\file{.prot}), the detailed protocol (\file{.log}) and all other output.
Copy it to start your own calculation. VESPAW reports energies in eV,
lengths in \AA, forces in eV/\AA{} and stress in GPa.

The checker also compares the results with the values quoted in this
tutorial and reports \texttt{PASS} or \texttt{FAIL}. If the code changes
so that an example no longer gives the numbers printed here, the checker
says so. The numbers in this tutorial were obtained on an Apple M-series
laptop (8 threads); timings are indicative.

The theory behind all examples is in the companion book
\emph{VESPAW: Theory and Methods} (\file{doc/theory.pdf}); every
input keyword is described in the input manual \file{doc/manual.pdf}.

% =====================================================================
\chapter{Installation}
% =====================================================================

\section{One-step installation}

\begin{shell}
git clone https://github.com/m-sotoudeh/VESPAW.git && cd VESPAW
./install.sh
export PATH="$HOME/.vespaw/bin:$PATH"
vespaw doctor
\end{shell}
\file{install.sh} finds or installs what VESPAW needs (Fortran compiler,
BLAS/LAPACK, FFTW, HDF5, MPI, Python) without administrator rights: on
macOS with Homebrew it uses Homebrew, otherwise a private conda-forge
environment in \file{\textasciitilde/.vespaw}; on clusters, load your
modules and run \texttt{./install.sh --system}. It builds the serial and,
if MPI works, the parallel version, and runs a short check.
\texttt{vespaw doctor} reports the state of the installation.

\section{Building by hand}

The build uses a parameter file for the toolchain, in \file{arch/}:
\begin{shell}
./configure --with-parmfile=arch/parms.arm64     # Apple Silicon
make fast_new             # -> bin/<arch>/vespaw_fast.x
make fast_parallel_new    # -> bin/<arch>/vespaw_fast_mpi.x (MPI)
make tools scripts        # analysis tools and the vespaw dispatcher
\end{shell}
Put \file{bin/<arch>/} on your \texttt{PATH}. Options and known build
problems are listed in \file{doc/INSTALL.md}.

\section{Testing the build}

\begin{shell}
python3 testsuite/run.py                  # the full test suite
python3 doc/src/tutorial/check_examples.py 02_si:scf
\end{shell}
The test suite runs every optimizer on a set of systems and checks
energies, gaps, forces and stress against reference values and against
each other; all cases should report \texttt{PASS}.

\section{Running}

A calculation needs a control file (\file{.cntl}) and a structure file
(\file{.strc}) with the same stem:
\begin{shell}
vespaw si.cntl                              # serial, OpenMP threads
OMP_NUM_THREADS=8 vespaw si.cntl
OPENBLAS_NUM_THREADS=1 mpirun -np 4 vespaw si.cntl   # MPI
\end{shell}

% =====================================================================
\chapter{Input basics}
% =====================================================================

\section{Two files, blocks and keywords}

Input files consist of nested blocks. A block starts with
\texttt{!NAME} and ends with \texttt{!END}; the file ends with
\texttt{!EOB}. Keywords are written \texttt{KEY=value}; strings are in
single quotes, logicals are \texttt{T} or \texttt{F}.

\paragraph{The control file} (\file{.cntl}) says \emph{what to do}:
\begin{itemize}
\item \key{!GENERIC}: number of steps \key{NSTEP}, time step \key{DT},
  \key{START=T} (start from scratch) or \key{F} (continue from the
  restart file), \key{FORCES=T}, \key{STRESS=T} (print them at the end).
\item \key{!FOURIER}: plane-wave cutoff \key{EPWPSI} in Rydberg and
  \key{CDUAL} (density cutoff = \key{CDUAL}$\times$\key{EPWPSI}).
\item \key{!DFT}: the functional, by name or number, e.g.\
  \key{TYPE='PBE'} or \key{TYPE='PBEsol'} (the VASP tags \key{'PE'},
  \key{'PS'}, \dots\ also work; the list is in the manual); the
  local-hybrid block \key{!NTBO} inside it.
\item \key{!PSIDYN}: how the electrons are optimized
  (\key{OPTIMIZER}, \key{MIXER}, convergence criteria).
\item \key{!MERMIN}: electronic temperature (metals).
\item \key{!RDYN}, \key{!CELL}: motion of atoms and cell.
\end{itemize}

\paragraph{The structure file} (\file{.strc}) says \emph{what the
system is}:
\begin{itemize}
\item \key{!GENERIC LUNIT[AA]}: length unit of all coordinates, here
  \AA{}ngstr\"om (\key{LUNIT[AA]=1.0}).
\item \key{!LATTICE T=}: the three lattice vectors (columns of
  $T$, in units of \key{LUNIT}).
\item \key{!KPOINTS DIV=}: the $k$-point mesh; add \key{SYMMETRY=T} to
  reduce it to its irreducible wedge using the space group.
\item \key{!OCCUPATIONS}: \key{NSPIN}, number of empty bands \key{EMPTY},
  total spin \key{SPIN[HBAR]}, charge.
\item \key{!SPECIES}: one block per element with the PAW setup
  (\key{!AUGMENT}) and, for hybrids, the local-orbital settings
  (\key{!NTBO}).
\item \key{!ATOM NAME= R=}: atoms and positions.
\end{itemize}

\paragraph{Switching blocks off.} Renaming a block (e.g.\
\texttt{!ORBPOT} $\to$ \texttt{!ORBPOT\_OFF}) makes the program ignore
it. \texttt{READ=F} inside a block also switches it off; this is easy to
do by accident.

\section{Choosing setups}

The directory \file{potentials} contains a species block for every
element. Copy the block of each element into the structure file.
\file{potentials/AUDIT.md} lists which setups passed the quality checks,
which are replaced, and how accurate they are in solids. Read it before
choosing a setup for a production calculation. When a setup is built,
the program prints a ghost test and a suggested cutoff (search the output
for \texttt{GHOST TEST} and \texttt{SUGGESTED ENMAX}).

\section{Convergence}

Two parameters must always be converged: the cutoff \key{EPWPSI} and
the $k$-point mesh. Converge the quantity you are interested in, not
only the total energy. Example~1 shows that forces can require a much
higher cutoff than the energy suggests, and Example~3 shows how the
energy of a metal oscillates with the $k$-point mesh.

% =====================================================================
\chapter{Example 1: a molecule (H\texorpdfstring{$_2$}{2}O)}
% =====================================================================

\textbf{Directory:} \file{examples/01\_h2o}. \textbf{You learn:} setting
up a molecule in a box, converging the cutoff with forces, relaxing a
geometry.

\section{Input}

The molecule sits in a cubic box of 10~\AA{} with $\Gamma$-point
sampling. It starts from a distorted geometry (O--H 1.00~\AA, angle
110$^\circ$).
\deck{../../../examples/tutorial/01_h2o/h2o.strc}
The control file:
\deck{../../../examples/tutorial/01_h2o/h2o.cntl}
The run \texttt{scf} uses Davidson with the variational optimizer:
\begin{verbatim}
!PSIDYN OPTIMIZER='DAV' MIXER='VARIATIONAL' MIXG0SQ=0.05 SAFEORTHO=T STOP=T
\end{verbatim}
\key{MIXG0SQ} sets the Kerker preconditioner. Its default (1.5) is made
for solids; for molecules and slabs with a lot of vacuum use a small
value (0.05--0.1), otherwise the variational optimizer stalls (see
\file{doc/KNOWN\_ISSUES.md}). Pulay mixing (\key{MIXER='PULAY'}) is a
faster alternative for molecules.

\section{Results}

\begin{shell}
python3 doc/src/tutorial/check_examples.py 01_h2o
    scf: E = -471.41787880 eV, 249 it, gap 5.8369 eV
    relax: E = -471.48491560 eV, 446 it, gap 5.9828 eV
\end{shell}
The relaxed geometry (in \file{work/relax/h2o.strc\_out}) has O--H
$=0.971$~\AA{} and H--O--H $=104.6^\circ$, in agreement with PBE
(0.970~\AA, 104.2$^\circ$).

\paragraph{The cutoff lesson.} At 30~Ry the same relaxation gives
O--H $=1.007$~\AA{} and $102.0^\circ$, a bond 4\% too long. At the PBE
geometry the forces are 2.6~eV/\AA{} at 30~Ry and 0.05~eV/\AA{} at 60~Ry.
Oxygen needs a high cutoff; converge it with forces, not only energies.

\section{The relaxation}

The run \texttt{relax} has two stages. Stage 1 converges the electrons
(as \texttt{scf}); stage 2 continues from the restart file
(\key{START=F}) and relaxes the atoms by damped dynamics of wave
functions and atoms together:
\begin{verbatim}
!PSIDYN STOP=T FRIC=0.02 SAFEORTHO=T
  !AUTO FRIC(-)=0.3 FACT(-)=0.97 FRIC(+)=0.3 FACT(+)=1. MINFRIC=0.02 !END
!END
!RDYN STOP=T FRIC=0.05
  !AUTO FRIC(-)=0.3 FACT(-)=0.97 FRIC(+)=0.3 FACT(+)=1. MINFRIC=0.05 !END
!END
\end{verbatim}
The \key{!AUTO} blocks (``autopilot'') lower the friction while the
energy decreases and stop the run when it has converged.
\key{FORCES=T} prints the final forces (in eV/\AA, also in section 07 of
the \file{.prot} report):
\begin{verbatim}
FINAL FORCES [EV/AA]
FORCE O_1   ...
\end{verbatim}
They are below 0.05~eV/\AA{} at the end of the relaxation.

% =====================================================================
\chapter{Example 2: a semiconductor (Si)}
% =====================================================================

\textbf{Directory:} \file{examples/02\_si}. \textbf{You learn:} total
energy, band gap, equation of state, forces and stress.

\section{Input}
\deck{../../../examples/tutorial/02_si/si.strc}
\deck{../../../examples/tutorial/02_si/si.cntl}
The primitive cell of diamond Si has two atoms. The placeholders are the
cutoff (\texttt{@EPW@}), the lattice scale (\texttt{@LUNIT@}), the
position of the second atom (\texttt{@X2@}) and extra output
(\texttt{@OUTPUT@}).

\section{Energy and band gap}

\begin{shell}
python3 doc/src/tutorial/check_examples.py 02_si:scf
    scf: E = -216.95911283 eV, 13 it, gap 0.7097 eV
\end{shell}
The variational optimizer converges in 13 iterations. The Kohn--Sham gap
of 0.71~eV on the $6\times6\times6$ mesh is the smallest gap between the
$k$-points of the mesh (the true PBE gap of Si is about 0.6~eV, between
$\Gamma$ and a point near X that the mesh does not contain). Section 05b
of the \file{.prot} report lists the band energies and occupations at
every $k$-point.

\section{Equation of state}

The runs \texttt{eos\_0.97} \dots\ \texttt{eos\_1.03} scale the lattice
by \texttt{@LUNIT@} at 40~Ry:
\begin{center}
\begin{tabular}{cc}\toprule
$a$ [\AA] & $E$ [eV] \\ \midrule
5.305 & $-216.86290$ \\
5.387 & $-216.94416$ \\
5.469 & $-216.97061$ \\
5.551 & $-216.94984$ \\
5.633 & $-216.88874$ \\ \bottomrule
\end{tabular}
\end{center}
A cubic fit of $E(V)$ gives $a_0=5.472$~\AA{} and $B=93$~GPa (all-electron
PBE: 5.469~\AA, 89~GPa; experiment: 5.431~\AA). A Birch--Murnaghan fit
is available with \file{vespaw\_murnaghan.x}. Energy--volume curves
need a converged cutoff: at fixed cutoff the number of plane waves
changes with the volume (Pulay stress).

\section{Forces and stress}

The run \texttt{forces} displaces the second atom by 0.05~\AA{} along $x$
and sets \key{FORCES=T STRESS=T}:
\begin{verbatim}
FORCE SI1   0.6907426272   -0.0000290266   -0.0000008242
FORCE SI2  -0.6907157157    0.0000335846   -0.0000830316
STRESS      2.1972555    -0.0003638     0.0001797   ...
STRESS     -0.0003627     2.3652870    -2.0113901   ...
STRESS      0.0001752    -2.0113887     2.3651746   ...
PRESSURE[GPA]      2.3092391  VOLUME[AA^3]      40.02575175
\end{verbatim}
The forces (eV/\AA) are equal and opposite, as they must be. The stress
is $\sigma=-(1/V)\,dE/d\epsilon$ in GPa, positive when the cell wants to
expand. The displacement along $x$ also produces a shear stress
$\sigma_{yz}$: in diamond, an internal displacement couples to the shear
strain. The test case \file{testsuite/si\_forces} checks forces
and stress against finite differences of the energy.

% =====================================================================
\chapter{Example 3: a metal (Al)}
% =====================================================================

\textbf{Directory:} \file{examples/03\_al}. \textbf{You learn:} electronic
temperature and $k$-point convergence of a metal.

\deck{../../../examples/tutorial/03_al/al.strc}
\deck{../../../examples/tutorial/03_al/al.cntl}
A metal has partially filled bands. \key{!MERMIN T[K]=1000.} gives the
electrons a temperature, i.e.\ Fermi occupations; the program minimizes
the free energy $A=E-TS$. The runs vary the $k$-point mesh:
\begin{center}
\begin{tabular}{ccc}\toprule
mesh & $A$ [eV], $T=1000$~K & \\ \midrule
$4^3$ & $-57.02980$ & \\
$6^3$ & $-57.36561$ & \\
$8^3$ & $-57.29449$ & \\
$12^3$ & $-57.30689$ & \\
$12^3$, $T=300$~K & $-57.30277$ & \\ \bottomrule
\end{tabular}
\end{center}
The energy does not converge monotonically: the $6^3$ mesh happens to be
59~meV too low. Only the difference between successive dense meshes tells
you whether you are converged ($8^3\to12^3$: 12~meV). A lower
temperature needs a denser mesh.

% =====================================================================
\chapter{Example 4: a magnet (NiO)}
% =====================================================================

\textbf{Directory:} \file{examples/04\_nio}. \textbf{You learn:}
antiferromagnetic order, seeding the spin state, complete setups.

\deck{../../../examples/tutorial/04_nio/nio.strc}
NiO orders antiferromagnetically: ferromagnetic (111) planes of Ni with
alternating spin. The cell contains two Ni atoms, \texttt{NI1} (spin up)
and \texttt{NI2} (spin down). The block
\begin{verbatim}
!ORBPOT
  !POT  ATOM='NI1' VALUE=+0.1 TYPE='D' S=1 RC=2. !END
  !POT  ATOM='NI2' VALUE=-0.1 TYPE='D' S=1 RC=2. !END
!END
\end{verbatim}
adds a small spin-dependent potential on the Ni $d$ orbitals that pushes
the system into the desired spin order. It is used only in the first
stage and then switched off by renaming it (\texttt{!ORBPOT\_OFF}); the
placeholder \texttt{@ORBPOT@} does exactly this.

\begin{shell}
python3 doc/src/tutorial/check_examples.py 04_nio
    dav_var: E = -3030.79719315 eV, 274 it, gap 1.4979 eV
    dav_coldseed: E = -3030.79713862 eV, 82 it, gap 1.5002 eV
\end{shell}
\texttt{dav\_var} seeds the spin order with friction dynamics and the
orbital potential, then converges with Davidson; \texttt{dav\_coldseed}
seeds and converges with Davidson directly. Both reach the same state.
The \file{.prot} report lists the magnetic moments: about
$\pm1.3~\mu_B$ in the Ni spheres, zero in total. The PBE gap of 1.5~eV is
far below the experimental 4.3~eV; see
Example~5 for hybrids.

\paragraph{Complete setups.} The Ni setup has two partial waves per
angular momentum (\key{NPRO=2 2 2}). With a single $s$ projector the
calculation can collapse into an unphysical state far below the ground
state; the setup generator warns about this (\texttt{INCOMPLETE SETUP}).

% =====================================================================
\chapter{Example 5: hybrid functionals (MgO)}
% =====================================================================

\textbf{Directory:} \file{examples/05\_mgo\_hybrid}. \textbf{You learn:}
the local hybrids \texttt{PBE0r} and \texttt{PBE-SL0r}.

\deck{../../../examples/tutorial/05_mgo_hybrid/mgo.cntl}
The placeholder \texttt{@NTBO@} is empty for PBE, and for the hybrids
\begin{verbatim}
# PBE0r
!NTBO MODUS='HYBRID' HFWEIGHT=0.25 K2=-0.25 SCALERCUT=2. !END
# PBE-SL0r
!NTBO MODUS='HYBRID' HFWEIGHT=0.25 DCTYPE='SCANL' K2=-0.25 SCALERCUT=2. !END
\end{verbatim}
\key{HFWEIGHT} is the fraction of exact exchange. \texttt{PBE0r} replaces
part of the onsite PBE exchange of the local orbitals by exact exchange;
\texttt{PBE-SL0r} in addition uses SCAN-L instead of PBE for that onsite
exchange (\key{DCTYPE='SCANL'}), with PBE everywhere else. A third local
hybrid, \texttt{SCAN0r}, uses SCAN everywhere (\key{TYPE='SCAN'
TSCANDELOC=T}); it is still experimental. The local orbitals are
defined per species in the structure file:
\deck{../../../examples/tutorial/05_mgo_hybrid/mgo.strc}
\key{NOFL} gives the number of local orbitals per angular momentum,
\key{LHFWEIGHT} a species-specific exchange fraction, \key{CV=T} includes
core--valence exchange.
\begin{center}
\begin{tabular}{lc}\toprule
functional & gap [eV] \\ \midrule
PBE & 4.326 \\ PBE0r & 4.034 \\ PBE-SL0r & 4.116 \\ \bottomrule
\end{tabular}
\end{center}
The local hybrids act only on the orbitals of the atoms (onsite). For
MgO, a closed-shell oxide, the correction is small; for oxides with open
$d$ shells (NiO, manganites) it is large. Take band gaps of hybrid
calculations from Davidson, not from RMM-DIIS, which can miss the lowest
empty band with hybrids.

% =====================================================================
\chapter{Example 6: molecular dynamics (Si\texorpdfstring{$_8$}{8})}
% =====================================================================

\textbf{Directory:} \file{examples/06\_md}. \textbf{You learn:}
Car--Parrinello molecular dynamics with a thermostat, and a pitfall of
small cells.

\deck{../../../examples/tutorial/06_md/si8.cntl}
Stage 1 converges the electrons. Stage 2 runs Car--Parrinello dynamics
without friction on the wave functions (\key{FRIC=0.}), with time step
\key{DT=10.0}~a.u.\ (0.24~fs), random initial velocities for 2000~K
(\key{RANDOM[K]}; equipartition quickly halves the temperature) and a
Nos\'e thermostat at 1000~K:
\begin{verbatim}
!RDYN STOP=T RANDOM[K]=2000.
  !THERMOSTAT T[K]=1000. FREQ[THZ]=10. !END
!END
\end{verbatim}
The detailed protocol (\file{.log}) prints a line per step, in atomic
units (Hartree):
\begin{verbatim}
    NFI  T[PSEC] T[K]  EKIN(PSI)       E(RHO)        ECONS   ANNEE   ANNER
!>   100   0.02419   696   0.00229    -31.82720    -31.77515 0.00000 0.00143
!>   500   0.12094  1160   0.00394    -31.83888    -31.77518 0.00000 0.00307
\end{verbatim}
with time, ionic temperature, fictitious kinetic energy of the wave
functions \texttt{EKIN(PSI)}, the potential energy \texttt{E(RHO)} and the
conserved energy \texttt{ECONS}. \texttt{ECONS} must stay constant; here it
changes by $3\times10^{-5}$~H (0.8~meV) in 400 steps. \texttt{EKIN(PSI)} must stay
small compared with the ionic kinetic energy (adiabaticity).

\paragraph{The small-cell pitfall.} With $\Gamma$-point sampling only,
the same run lets the energy fall by 6.8~eV below that of diamond: in an
8-atom cell sampled at $\Gamma$ alone, distorted structures are lower in
energy than diamond. With a $2\times2\times2$ mesh diamond is correctly
lower by 4.4~eV. For MD, use a cell large enough for $\Gamma$-point
sampling, or a $k$-point mesh.

% =====================================================================
\chapter{Choosing the optimizer and troubleshooting}
% =====================================================================

\section{Which optimizer?}

\begin{center}
\begin{tabular}{p{3.2cm}p{11cm}}\toprule
\key{OPTIMIZER='DAV'} & block Davidson. Robust from any start; the
  reference. Default choice for small and medium systems. \\
\key{OPTIMIZER='RMM'} & RMM-DIIS with a Davidson warm-up. Fastest for
  large systems (64 atoms and more). Not for band gaps with hybrids. \\
\key{OPTIMIZER='CG'} & conjugate gradients, band by band. Reliable,
  slower for many bands. \\
dynamics (default) & damped Car--Parrinello dynamics
  (\key{FRIC}, \key{!AUTO}). Slow but very robust; required for
  molecular dynamics and convenient for relaxations. \\
\key{OPTIMIZER='EDFT'} & direct joint minimization of orbitals and
  occupations (Route~2). Experimental. \\ \bottomrule
\end{tabular}
\end{center}
With the eigensolvers the mixer is \key{MIXER='VARIATIONAL'} (default):
every accepted energy is an upper bound of the ground-state energy and
decreases monotonically. \key{MIXER='PULAY'} and \key{'ANDERSON'} are
classical density mixing. Convergence is controlled by \key{EDIFF}
(energy change, default $2.7\times10^{-5}$~eV) and \key{MIXTOL} (residual).

\section{Troubleshooting}

\paragraph{The SCF does not converge.}
\begin{itemize}
\item Molecules and slabs: reduce \key{MIXG0SQ} to 0.05--0.1.
\item Metals: set an electronic temperature (\key{!MERMIN}).
\item Magnetic systems: seed the spin order (\key{!ORBPOT}, Example~4),
  or start with a few steps of damped dynamics.
\item Compare with damped dynamics. If the variational optimizer finds
  an energy \emph{below} the dynamics result without converging, the
  potential is not the derivative of the energy; that is a bug worth
  reporting.
\end{itemize}

\paragraph{The energy is far too low.} A collapse into an unphysical
state: check the setup output for \texttt{INCOMPLETE SETUP} or
\texttt{GHOST}, use complete setups (two projectors per angular
momentum), and look for very deep eigenvalues.

\paragraph{Wrong structure.} Converge the cutoff with forces
(Example~1) and the $k$-points (Example~3); for $3d$ metals read the
accuracy notes in \file{potentials/AUDIT.md}.

\paragraph{Orthogonalization fails at the start.}
\texttt{LOOP FOR ORTHOGONALIZATION IS NOT CONVERGED} in damped dynamics
from random wave functions: start with a few Davidson steps, or with
\key{SAFEORTHO=F}.

% =====================================================================
\chapter{Output and analysis}
% =====================================================================

\section{Files}

\begin{center}
\begin{tabular}{ll}\toprule
\file{.prot} & report: input, system, model, SCF history, energies,
  $k$-points and occupations, \\ & atomic charges and moments, forces, stress, timing \\
\file{.log} & detailed protocol: setup construction, per-step reports \\
standard output & live table of iterations and the final energy \\
\file{results.h5} & numerical results (with \texttt{!HDF5OUT}) \\
\file{.rstrt} & restart file (wave functions, positions) for \key{START=F} \\
\file{.strc\_out} & final structure \\
\file{.pdos} & projections for the density of states \\
\file{.tra} & trajectories (MD, relaxation) \\ \bottomrule
\end{tabular}
\end{center}

\section{Density of states}

The DOS tool reads the \file{.pdos} file. For Si
(\file{examples/02\_si/si.dcntl}):
\deck{../../../examples/tutorial/02_si/si.dcntl}
\begin{shell}
cd examples/tutorial/02_si/work/scf
vespaw_dos.x ../../si.dcntl
\end{shell}
writes \file{dos\_total.dos}, \file{dos\_si\_s.dos} and
\file{dos\_si\_p.dos} (energy in eV, DOS, integrated DOS). Weights can be
atoms and angular momenta (\key{TYPE='S','P','D'}) or single orbitals
(\key{!ORB}, e.g.\ $e_g$ and $t_{2g}$); see \file{doc/manual.pdf},
``The Density of States Analysis Tool''.

\section{HDF5 and Python}

With \texttt{!HDF5OUT READ=T FILE='results.h5' !END} in the
\key{!CONTROL} block the program writes energy, forces, eigenvalues and
provenance to an HDF5 file. The Python package reads it:
\begin{shell}
pip install -e python/
export PATH="$HOME/.vespaw/bin:$PATH"     # after install.sh
\end{shell}
\begin{verbatim}
from vespaw import Calculation
r = Calculation("si.cntl").run()
print(r.energy, r.gap, r.forces)   # eV, eV, eV/Angstrom
\end{verbatim}
For the Si example this prints $E=-216.95911$~eV and a gap of
0.7097~eV. A YAML input front end is described in
\file{doc/USAGE.md}.

\end{document}
